\documentclass[12pt]{book}

\usepackage{amsmath, amssymb, amsthm, mathtools}
\usepackage{tcolorbox}
\tcbuselibrary{theorems,breakable}
\usepackage{xcolor}

\newcommand{\R}{\mathbb{R}}
\newcommand{\N}{\mathbb{N}}
\newcommand{\Z}{\mathbb{Z}}
\newcommand{\Q}{\mathbb{Q}}

\definecolor{theoremcolor}{RGB}{0,102,204}      % Blue
\definecolor{definitioncolor}{RGB}{0,153,76}    % Green
\definecolor{lemmacolor}{RGB}{153,0,153}        % Purple
\definecolor{explanationcolor}{RGB}{102,102,102} % Gray

\newtcbtheorem[number within=section]{theorem}{Theorem}%
{colback=theoremcolor!5!white,
 colframe=theoremcolor!75!black,
 colbacktitle=theoremcolor!85!black,
 fonttitle=\bfseries,
 boxrule=1pt,
 breakable
}{th}

\newtcbtheorem[number within=section]{definition}{Definition}%
{colback=definitioncolor!5!white,
 colframe=definitioncolor!75!black,
 colbacktitle=definitioncolor!85!black,
 fonttitle=\bfseries,
 boxrule=1pt,
 breakable
}{def}

\newtcbtheorem[number within=section]{lemma}{Lemma}%
{colback=lemmacolor!5!white,
 colframe=lemmacolor!75!black,
 colbacktitle=lemmacolor!85!black,
 fonttitle=\bfseries,
 boxrule=1pt,
 breakable
}{lem}

\newtcolorbox{explanation}{
 colback=explanationcolor!5!white,
 colframe=explanationcolor!75!black,
 colbacktitle=explanationcolor!85!black,
 title={Explanation},
 fonttitle=\bfseries,
 boxrule=1pt,
 breakable
}

\newcommand{\btheorem}{\begin{theorem}{}{}~}
\newcommand{\etheorem}{\end{theorem}}
\newcommand{\bdefn}{\begin{definition}{}{}~}
\newcommand{\edefn}{\end{definition}}
\newcommand{\blemma}{\begin{lemma}{}{}~}
\newcommand{\elemma}{\end{lemma}}
\newcommand{\bexpl}{\begin{explanation}~}
\newcommand{\eexpl}{\end{explanation}}

% Advanced versions with title and label
\newcommand{\btheoremtl}[2]{\begin{theorem}{#1}{#2}~}
\newcommand{\bdefntl}[2]{\begin{definition}{#1}{#2}~}
\newcommand{\blemmatl}[2]{\begin{lemma}{#1}{#2}~}

% Book information
\title{Anvaya of Mathematical Definitions}
\author{Bikash Thakur}
\date{September 18 2025}

\begin{document}

% Title page
\maketitle

% Table of contents
\tableofcontents

\btheorem
$T:K^n \to K^m$. Then $R(T)$ is the span of the columns of the matrix.
\etheorem

\bobservation
If $C$ is compact then conv(C) is also compact.
\eobservation
\bremark
C = conv(C)
\eremark
\breason
Caratheodory
\ereason
\bdefn
what are extreme points? Pizza slice. A point $x$ in a convex set $C$ is called an extreme point if it has the following property \begin{enumerate} \item $x = \lambda y + (1-\lambda)z, yz \in C, \lambda \in (0, 1)$ then $y=z=x$. Open convex sets such as an open circle have no extreme points. Half-space has no extreme points as it is not compact.



\edefn

Faces are generalizations of extreme points. 

\bdefn
		A set $F\subset C$ of a convex set $C$ is called a face of $C$ if for any $x \in F$ the following holds \[x = \lambda y + (1-\lambda) z\] for $y,z \in C$ and $\lambda \in (0, 1)$. Then $y,z \in F$. Extreme points are faces that have $1$ element. Faces are like maximal. A triangle with three corners extreme points. Example of the krein-milman theorem.

\edefn

\btheorem
		Milman's Theorem (Finite Dimensional)
		Let $C$ (not necessarily convex set) be closed and bounded (compact) set in $\R^d$. If $x$ is an extreme point of conv(C) then $x \in C$.

		\etheorem

\begin{theorem}
	Krein-Milman (Finite dimensional version)
	Let $C$ be a compact, convex set in $\R^d$. Let $S$ be the set of extreme points of $C$. Then $S$ is non-empty, and $\conv(S) = C.$


	In Words: Closed convex sets ARE just convex hull of their extreme points. The true power of Krein-Milman comes --simplex(gen. of higher dim triangle). Triangles have a unique property that every point is a unique convex combination of the three extreme points.$
\end{theorem}


\begin{theorem}
	Projection Theorem: Let $C$ be a closed convex set and $y \notin C \subset \R^d$. Then $\exits$ a unique $!$ point $x^* \in C$ such that $x^*$ solves the following problem
	$\minimize \norm{x -y} \subjectto x \in C$.
	i.e.,
\end{theorem}

\bremark
Unfortunately if $C$ is compact, S is the set of extreme points of C, $S$ may not be compact.
Two ice-cream cones.
a single point on the straig line along the circle is not an extreme point.



		\begin{definition}
			If $X$ is a finite dimensional vector spaces with basis $B=\{x_1, \cdots, x_n\}$ then $B^* = \{x_1^*, \cdots, x_n^*\} \subset X'$ is the dual basis if \[x_i^*(x_j) = \delta_{ij}\]

		\end{definition}
		input is column and output is row; for operators T to create a matrix.
all rank 1 matrices can be written as a column times row vector.

		\begin{theorem}
			$T:X\to Y$ is a linear transform, $X,Y$ are finite dimensional vector spaces. Let $B_X = \{x_1, \cdots, x_n\}$ be a basis of $X$ with the dual $B_X^* = \{x_1^*, \cdots, x_n^*\}$. Let $B_Y=\{y_1, \cdots, y_n\}$ be a basis of $Y$ with the dual basis $B_Y^* = \{y_1^*, \cdots, y_n^*\}$. Let $K^n$ have standard basis, $\{e_1, \cdots, e_n\}$ and their duals $\{e_1^*, \cdots, e_n^*\}$. Define $J_X x_i = e_i, J_Y y_i = \tilde{e_i}, J_X: X \to K^n , J_Y: Y \to K^m\}$. Then
			\begin{enumerate}
				\item \[T = \sum_{i=1}^m \sum_{j=1}^n T_{ij} y_i x_j^*\]
				\item $\tilde{T}: J_Y T J_X^{-1}: K^n \to K^m$ then $\tilde{T} = \sum_{i=1}^m \sum_{j=1}^n T_{ij} \tilde{e_i} e_j^*$
					$T_{ij} = y_i^* (Tx_j)$ by def.


			\end{enumerate}

		\end{theorem}

		Recall: $T:X \to Y$ X, Y are finited dimensional. $T': Y' \to X'$. $(T'l)': X \to K$ and $(T'l)(x) = (l, Tx)$ so $(T'l, x) = (l, Tx) = l(Tx)$.
		$X'' \sim X$, $\tilde{x} \in X''$ so $\tilde{x}(l) = l(x)$.
		$Y'$ has a basis $B_Y^*$
		$Y''$ has a basis $B_Y$
		$X'$ has a basis $B_X^*$
		$X''$ has a basis $B_X$.
		$T_{ij}' = (T'y_j^*x_i) = (T'y_j^*) (x_i) = (T'y_j^*, x_i) = (y_j^*, Tx_i) = T_{ji}.$

		$T: K^n \to K^m$, $T^T: K^m \to K^n$, $(T^T)_{ij} = T_{ji}$
		$T': Y' \to X', J_Y': Y' \to K^m, J_X' : X' \to K^n$ and $



\end{document}
